\begin{proposition}\label{sga3:XXI:7.1.3}
Let \(\cR\) be a root datum. The following conditions are
equivalent:
\begin{enumerate}
\item[\SGAInlineRefTarget{sga3:XXI:7.1.3:item:i}(i)] There is no nontrivial partition of \(R\) into two
  closed symmetric subsets.
\item[\SGAInlineRefTarget{sga3:XXI:7.1:item:ii}(ii)] For some (respectively every) system of simple roots
  \(\Delta\) of \(R\), there is no partition of \(\Delta\) into
  two nonempty orthogonal subsets.
\item[\SGAInlineRefTarget{sga3:XXI:7.1.3:item:iii}(iii)] The natural representation of \(W(\cR)\) on
  \(\VR(R)\) is irreducible.
\item[\SGAInlineRefTarget{sga3:XXI:7.1.3:item:iv}(iv)] For every pair \((\alpha,\beta)\) of roots, there is a
  sequence of roots
  \[
   \alpha_0,\alpha_1,\ldots,\alpha_n,
   \qquad
   \alpha_0=\alpha,\quad\alpha_n=\beta,
  \]
  such that \(\alpha_i\) and \(\alpha_{i+1}\) are nonorthogonal
  for \(i=0,\ldots,n-1\).
\end{enumerate}
\end{proposition}

\begin{proof}
The equivalence \SGARefLink{sga3:XXI:7.1.3:item:i}{(i)}\(\Leftrightarrow\)\SGARefLink{sga3:XXI:7.1:item:ii}{(ii)} follows from
\ref{sga3:XXI:7.1.1}, and \SGARefLink{sga3:XXI:7.1.3:item:iv}{(iv)}\(\Rightarrow\)\SGARefLink{sga3:XXI:7.1.1:item:ii}{(ii)} is immediate.
Conversely, if \SGARefLink{sga3:XXI:7.1.1:item:ii}{(ii)} holds for \(\Delta\), then \SGARefLink{sga3:XXI:7.1.3:item:iv}{(iv)} holds when
\(\alpha,\beta\in\Delta\). Every root is nonorthogonal to some
simple root, by \ref{sga3:XXI:3.1.1}, for example. Moreover,
\SGARefLink{sga3:XXI:7.1.3:item:iii}{(iii)}\(\Rightarrow\)\SGARefLink{sga3:XXI:7.1.3:item:i}{(i)}: under the hypotheses of
\ref{sga3:XXI:7.1.1}, the space \(\VR(R')\) is stable under
\(W(\cR)\). It remains to prove \SGARefLink{sga3:XXI:7.1.1:item:i}{(i)}\(\Rightarrow\)\SGARefLink{sga3:XXI:7.1.3:item:iii}{(iii)}.

Let \(H\) be a vector subspace of \(\VR(R)\), stable under
\(W(\cR)\). For every \(\alpha\in R\), the equality
\(s_\alpha(H)=H\) immediately gives
\[
 \alpha\in H
 \qquad\text{or}\qquad
 \alpha^\ast\in H^\perp,
\]
where \(H^\perp\) denotes the orthogonal complement of \(H\) in
\(\VR(R^\ast)\), which is in duality with \(\VR(R)\). Put
\[
 R'=\{\alpha\in R\mid\alpha\in H\},
 \qquad
 R''=\{\alpha\in R\mid\alpha^\ast\in H^\perp\}.
\]
This gives a partition of \(R\) into two closed symmetric subsets.
\end{proof}

\begin{definition}\label{sga3:XXI:7.1.4}
A root datum (respectively root system) satisfying the equivalent
conditions of \ref{sga3:XXI:7.1.3} and having nonzero semisimple
rank is called \emph{irreducible}.
\end{definition}

\begin{corollary}\label{sga3:XXI:7.1.5}
For every root datum \(\cR\), there is a unique partition, up to
order, of \(R\) into closed symmetric irreducible subsets.
\end{corollary}

\begin{corollary}\label{sga3:XXI:7.1.6}
Every adjoint (respectively simply connected) root datum is a
product of irreducible adjoint (respectively simply connected) root
data.
\end{corollary}

\begin{proof}
It suffices to consider the adjoint case. The assertion then follows
from the fact that, under the hypotheses of
\ref{sga3:XXI:7.1.1},
\[
 \Gamma_0(R)=\Gamma_0(R')\times\Gamma_0(R'').
\]
\end{proof}

\begin{corollary}\label{sga3:XXI:7.1.7}
For every root datum (respectively reduced root datum) \(\cR\),
there is an isogeny
\[
 \cR\longrightarrow\cR',
\]
where \(\cR'\) is a product of a trivial root datum and irreducible
simply connected root data (respectively irreducible reduced simply
connected root data).
\end{corollary}

\subsection{Properties of irreducible root data}
\label{sga3:XXI:7.2}

\begin{definition}\label{sga3:XXI:7.2.1}
Let \(\cR\) be an irreducible root datum. For every \(\alpha\in R\),
put
\[
 \operatorname{long}(\alpha)
 =\ell(\alpha)/\ell(\alpha_0),
\]
where \(\alpha_0\in R\) is such that \(\ell(\alpha_0)\) is minimal.
The number \(\operatorname{long}(\alpha)\) is called the
\emph{length} of \(\alpha\).
\end{definition}

\begin{lemma}\label{sga3:XXI:7.2.2}
Let \(\cR\) be an irreducible root datum. The Weyl group acts
transitively on the set of roots of a given length.
\end{lemma}

\begin{proof}
Let \(\alpha,\beta\in R\). Since the representation of \(W(\cR)\)
on \(\VR(R)\) is irreducible, \(\alpha\) cannot be orthogonal to
all \(w(\beta)\), \(w\in W(\cR)\). Hence some \(w(\beta)\) is
nonorthogonal to \(\alpha\). Since
\[
 \ell(w(\beta))=\ell(\beta),
\]
the assertion follows from \ref{sga3:XXI:2.3.2}.
\end{proof}

\begin{lemma}\label{sga3:XXI:7.2.3}
If \(\cR\) is irreducible and reduced, then
\[
 \operatorname{long}(R)
 \in\bigl\{\{1\},\{1,2\},\{1,3\}\bigr\}.
\]
\end{lemma}

\begin{proof}
By the observation used above, for every
\(\alpha,\beta\in R\) there is \(w\in W(\cR)\) such that
\(w(\beta)\) is nonorthogonal to \(\alpha\). Therefore, by
\ref{sga3:XXI:2.3.1},
\[
 \frac{\ell(\alpha)}{\ell(\beta)}
 \in\left\{1,2,3,\frac12,\frac13\right\}.
\]
Thus \(\operatorname{long}(\alpha)\) is \(1\), \(2\), or \(3\).
But lengths \(2\) and \(3\) cannot both occur, since their ratio
\(2/3\) is not in the displayed set.
\end{proof}
